
\section{Pastor O. L. Anderson and the Guiding Principles of the Lutheran Free Church}

“Folkebladet,” July 17, 1918

Elias P. Harbo

Pastor L. O. Anderson has, in Folkebladet for June 26 of this year, given an account of why he left the Lutheran Free Church.

In this account he says, among other things: “In short: because I have come to realize that the manner in which I conduct congregational work is not in harmony with the view of congregational work maintained by the Guiding Principles of the Lutheran Free Church, but is rather in harmony with the view of congregational work maintained by the larger church bodies.”

The first thing we notice in the remarkable statement quoted above is that, in Pastor Anderson’s opinion, there exist two entirely different views of congregational work in “the larger church bodies” and in the Lutheran Free Church, in that “the larger church bodies” maintain one view and the Lutheran Free Church maintains another. And so entirely different are these two views of congregational work that Pastor Anderson could not remain in the Lutheran Free Church and at the same time conduct his congregational work according to the view of congregational work maintained in “the larger church bodies.” We are very grateful to Pastor Anderson for this information, namely, that he has made the discovery, or come to the recognition, that there exist two entirely different views of congregational work in “the larger church bodies” and in the Lutheran Free Church. And of course there is an entirely different practice as well, for it is surely right to conclude that each one practices according to his view of how congregational work ought to be carried on.

Pastor Anderson therefore does not agree with those who say that there is no difference between the Lutheran Free Church and the church bodies.

Pastor Anderson declares without any reservation that “the larger church bodies” and the Lutheran Free Church maintain two entirely different views of congregational work.

And it is this view of how work ought to be carried on for and in the congregation, in order that the work may bear fruit for the life and flourishing of the congregation, which is of such exceedingly great importance to all of us who are concerned for God’s congregation, its growth and flourishing.

But the view a man has of the work for and in the congregation depends, of course, upon the view he has of the congregation itself, what the congregation is according to its essential nature, and what its calling and task are.

Pastor Anderson gives no explanation of wherein this difference between the two views consists. He gives only the declaration that the Guiding Principles of the Lutheran Free Church maintain one view and “the larger church bodies” maintain another.

Nor does he mention the points in the Guiding Principles in which the Lutheran Free Church’s view is maintained.

We will therefore, through Folkebladet, address a friendly request to Pastor Anderson that he give a more detailed explanation of wherein these two different views consist, and also name the points in the Guiding Principles which, in his understanding, maintain the Lutheran Free Church’s view of the work for and in the congregation.

This question, of how we are to work for and in the congregation in such a way that it may appear among us more as a living and free congregation, appear in its New Testament form, appear Spirit-filled and vigorous in action, is a question which all of us who are born of God regard with warm interest, and to which we are therefore undoubtedly willing to contribute our share, so that all possible light may be shed upon it, in order that we may attain the greatest possible clarity in this matter, which is of such exceedingly great importance for the kingdom of God and the salvation of souls.

But this time we shall say no more about this question.